% Bab I Pendahuluan
% Template: TEMPLATE_BAB.md
% Gaya: WRITING_GUIDE.md
\cleardoublepage
\chapter{PENDAHULUAN}
\label{chap:pendahuluan}

Bab ini meletakkan dasar bagi keseluruhan laporan Tugas Akhir. Pembahasan dibuka dengan latar belakang yang memetakan tekanan praktik DataOps dan MLOps pada lingkungan produksi, kemudian diturunkan menjadi empat rumusan masalah dan empat tujuan yang saling berpasangan. Batasan masalah menetapkan lingkup pengerjaan, metodologi menjelaskan kerangka \textit{Design Science Research Methodology} yang diikuti, dan sistematika penulisan menutup bab dengan peta isi laporan. Penomoran RM-1 hingga RM-4, T-1 hingga T-4, serta batasan penelitian BP-1 hingga BP-3 yang diperkenalkan pada bab ini menjadi rujukan tetap bagi seluruh bab berikutnya, sedangkan batasan implementasi BI-1 dan BI-2 diperkenalkan bersama lingkungan implementasi pada Bab~\ref{chap:implementasi}.

\section{Latar Belakang}
\label{sec:latar-belakang}

Keberhasilan \textit{machine learning} pada lingkungan produksi ditentukan oleh kemampuan memindahkan model dari eksperimen ke layanan yang stabil, terpantau, dan dapat diperbarui, bukan oleh ketepatan model semata. \textcite{sculley2015hidden} menunjukkan kode model hanya sebagian kecil dari sistem nyata, sedangkan sisanya berupa pengelolaan data, konfigurasi, pemantauan, \textit{serving}, dan \textit{lineage}. \textcite{paleyes2022challenges} memetakan tantangan \textit{deployment} dari kualitas data hingga pemeliharaan model, sementara adaptasi DevOps oleh komunitas \textit{machine learning} melahirkan MLOps \autocite{kreuzberger2023mlops} dan DataOps tumbuh serupa untuk \textit{pipeline} data. \textcite{rella2022mlops} mencatat keduanya berbagi prinsip otomasi yang sama tetapi kerap berjalan sebagai dua \textit{tech stack} terpisah, sedangkan \textcite{jain2025integrating} menempatkan integrasinya sebagai syarat \textit{pipeline} yang skalabel dan tangguh. Upaya integrasi tersebut dihadapkan pada dua tekanan utama yang, ketika tidak teratasi, memicu tekanan ketiga berupa efek domino.

\vspace{0.5em}
\begin{enumerate}
    \item \textbf{Beban membangun platform dari nol di atas kumpulan komponen yang terfragmentasi}

    Kajian sistematis \textcite{najafabadi2024analysis} atas 43 studi primer mengategorikan 35 komponen arsitektur MLOps beserta ragam \textit{tool} pendukungnya, menandai luasnya ruang pilihan sebelum satu platform utuh terbentuk. \textcite{burgueno2025toward} menambahkan bahwa komponen \textit{open source} MLOps cenderung berdiri sendiri alih-alih membentuk kerangka kohesif, sehingga beban integrasi, kurva belajar, dan penyelarasan versi ditanggung tim pemakai sejak awal.

    \vspace{0.5em}

    \item \textbf{Biaya berlangganan layanan terkelola dan \textit{trade-off} antara waktu pengembangan dengan anggaran berjalan}

    Layanan terkelola mempercepat adopsi \autocite{li2023managed}, tetapi studi empiris \textcite{hamza2024serverlesscost} menemukan bahwa penagihan bayar per pakai dan pelaporan biaya \textit{cloud} yang tidak transparan menyulitkan estimasi anggaran dan dapat menjadi tidak hemat pada beban besar, sejalan dengan \textcite{adusumilli2025serverless}. Ditambah lagi dengan risiko keterikatan penyedia \autocite{mo2023serverless}, organisasi terjebak antara membangun tim platform dengan waktu pengembangan panjang atau berlangganan dengan biaya sulit diestimasi dan ruang kustomisasi yang sempit.

    \vspace{0.5em}

    \item \textbf{Efek domino fragmentasi yang merambat ke seluruh rantai DataOps dan MLOps}

    Ketika kedua tekanan tidak terjawab oleh satu \textit{control plane}, sambungan yang terputus antar \textit{layer} \textit{ingestion}, fitur, \textit{training}, dan \textit{serving} memaksa integrasi dirakit ulang tiap proyek, sementara metadata, identitas, dan \textit{lineage} dikelola sendiri-sendiri sehingga asal data sulit ditelusuri dan nilai fitur rawan berbeda antara \textit{training} dan \textit{serving}. \textcite{shankar2022operationalizing} mendokumentasikan lewat wawancara delapan belas insinyur bahwa pemeliharaan \textit{pipeline} dan integrasi antar komponen menjadi keluhan berulang, dengan rincian per \textit{layer} pada Bab~\ref{chap:analisis-masalah}.
\end{enumerate}

Gabungan ketiga tekanan itu menggerakkan pengembangan platform yang menyatukan DataOps dan MLOps di atas Kubernetes dengan komponen \textit{open source}. Kubernetes dipilih sebagai standar \textit{de facto} orkestrasi \textit{cloud-native} yang teradopsi luas di lingkungan \textit{enterprise} \autocite{lakka2025kubernetes}, sedangkan komponen \textit{open source} membuat arsitektur dapat ditiru tanpa keterikatan penyedia sekaligus menghapus biaya lisensi. Rancangan dijaga \textit{domain-agnostic} agar dapat diadaptasi lintas domain tanpa mengubah \textit{control plane}, dengan keluaran empat \textit{sub-sistem} pemutus rantai fragmentasi. Seluruh komponen ditata pada repositori Git internal dengan \textit{operator} GitOps sebagai mesin rekonsiliasi deklaratif, sehingga sembilan \textit{layer} arsitektur, dari orkestrasi hingga keamanan platform, terhubung pada satu jalur GitOps yang dirinci pada Bab~\ref{chap:perancangan}.


\section{Rumusan Masalah}
\label{sec:rumusan-masalah}

Latar belakang di atas mengarah pada empat rumusan masalah yang ingin diselesaikan oleh platform yang dikembangkan. Butir-butir tersebut dirumuskan sebagai pertanyaan agar setiap permasalahan memiliki batas teknis yang jelas dan rute jawaban berupa \textit{sub-sistem} atau kontrak antar komponen pada bab-bab berikutnya. Setiap pertanyaan menyasar properti sistem yang belum terpenuhi, bukan ketiadaan satu \textit{tool} tertentu, agar jawabannya berlaku lintas pilihan komponen dan menjaga sifat \textit{domain-agnostic} platform. Rumusan di bawah ini selanjutnya dirujuk dengan singkatan RM-1 hingga RM-4 pada bab-bab berikutnya.

\vspace{0.5em}
\begin{enumerate}
    \item Bagaimana menyatukan \textit{tech stack} DataOps dan MLOps yang terfragmentasi pada \textit{layer} \textit{ingestion}, \textit{processing}, \textit{storage}, \textit{training}, dan \textit{model serving} ke dalam satu \textit{control plane} pada Kubernetes?

    \item Bagaimana mewujudkan tata kelola data yang konsisten lintas siklus data, fitur, dan model sehingga katalog metadata, \textit{lineage}, dan kontrol kualitas tidak berjalan sebagai pekerjaan susulan?

    \item Bagaimana mengotomasi deteksi \textit{drift} dan menjamin \textit{point-in-time correctness} agar tidak terjadi kebocoran temporal antara mode \textit{batch} dan \textit{streaming}?

    \item Bagaimana menyajikan fitur multimodal, yang mencakup \textit{tabular feature}, \textit{aggregate feature}, dan \textit{vector embedding} berdimensi tinggi, dari satu sumber kebenaran dengan nilai yang konsisten antara fase \textit{training} dan \textit{serving}?
\end{enumerate}

Secara berurutan, keempat rumusan masalah di atas dijawab oleh empat tujuan pada Subbab~\ref{sec:tujuan}, dengan rincian kebutuhan fungsional dan nonfungsional diturunkan pada Bab~\ref{chap:analisis-masalah}. Rumusan-rumusan itu tidak saling tumpang tindih karena masing-masing menyasar satu titik kegagalan yang berbeda, yaitu perakitan lintas \textit{layer}, tata kelola data, deteksi \textit{drift} beserta pencegahan kebocoran temporal, dan \textit{feature serving} multimodal, sekaligus bersama-sama mencakup rentang dari siklus data hingga \textit{model serving}. Karena dirumuskan dengan batas teknis yang jelas dan berpijak pada kegagalan yang terdokumentasi pada latar belakang, keempat rumusan menjadi acuan tetap yang jawabannya ditunjukkan melalui artefak yang berjalan. Pemenuhan tiap rumusan ditelusuri kembali pada Bab~\ref{chap:evaluasi} melalui skenario uji yang menautkan rumusan dengan bukti verifikasinya.

\section{Tujuan}
\label{sec:tujuan}

Penelitian ini merumuskan empat tujuan yang masing-masing menjawab satu rumusan masalah, dengan keluaran berupa artefak teknis yang dapat ditunjukkan dan dipakai ulang. Setiap tujuan menetapkan keluaran berupa rancangan arsitektur, manifes \textit{deployment} pada repositori Gitea, dan kontrak antar komponen yang dapat diuji pada \textit{use case} verifikasi. Singkatan T-1 hingga T-4 dipakai untuk merujuk keempat tujuan di bawah ini pada bab-bab berikutnya, dengan pemetaan satu lawan satu terhadap RM-1 hingga RM-4. Tiap tujuan dirumuskan sebagai tindakan merancang dan mewujudkan artefak yang dapat ditunjukkan, bukan sebagai aktivitas mengukur atau menguji.

\vspace{0.5em}
\begin{enumerate}
    \item Merancang dan mewujudkan arsitektur platform DataOps dan MLOps terintegrasi di atas Kubernetes yang menyatukan \textit{data ingestion}, \textit{processing}, \textit{storage} dan \textit{feature store}, \textit{training}, \textit{model serving}, serta \textit{governance} dan \textit{observability} pada satu \textit{control plane} deklaratif dengan praktik GitOps, menjawab RM-1.

    \item Membangun \textit{sub-sistem} tata kelola data yang menyediakan katalog metadata, perekaman \textit{lineage}, dan kontrol kualitas data sebagai layanan bersama yang dijalankan secara otomatis di seluruh \textit{pipeline} data, fitur, dan model, menjawab RM-2.

    \item Mengimplementasikan \textit{sub-sistem} deteksi \textit{drift} terotomasi dengan kebijakan \textit{retraining} berbasis ambang batas serta menjamin \textit{point-in-time correctness} pada \textit{feature serving} lintas mode \textit{batch} dan \textit{streaming}, menjawab RM-3.

    \item Mengembangkan layanan fitur \textit{dual-store} yang melayani \textit{tabular feature} dan \textit{aggregate feature} melalui satu kontrak API dengan jaminan konsistensi nilai antara fase \textit{training} dan \textit{serving}, serta melayani \textit{vector embedding} berdimensi tinggi pada indeks vektor terpisah dengan ruang \textit{embedding} yang konsisten pada kedua fase tersebut, menjawab RM-4.
\end{enumerate}

Kriteria keberhasilan keseluruhan platform diukur dari empat sudut yang masing-masing menautkan satu tujuan. Setiap kriteria dirumuskan agar dapat diperiksa pada artefak yang berjalan, bukan pada rancangan di atas kertas. Pemeriksaannya dilakukan melalui skenario uji pada Bab~\ref{chap:evaluasi} dengan penerimaan fungsional dan struktural sesuai lingkup batasan penelitian. Kriteria-kriteria berikut menjadi tolok ukur minimum sebelum platform dinyatakan menjawab rumusan masalah.

\vspace{0.5em}
\begin{enumerate}
    \item Seluruh artefak dapat dipasang dari satu repositori Git melalui Argo CD tanpa intervensi manual setelah \textit{bootstrap} awal klaster (T-1).

    \item Seluruh \textit{pipeline} data, fitur, dan model terdaftar pada katalog metadata dengan \textit{lineage} yang dapat ditelusuri (T-2).

    \item Mekanisme deteksi \textit{drift} memicu \textit{retraining} secara otomatis ketika ambang batas dilewati pada \textit{use case} verifikasi (T-3).

    \item \textit{Tabular feature} dan \textit{aggregate feature} tersaji melalui satu kontrak API dengan nilai yang konsisten antara penyimpanan \textit{offline} dan \textit{online}, sedangkan \textit{embedding} tersaji pada indeks vektor terpisah dengan ruang \textit{embedding} yang konsisten antara fase \textit{training} dan \textit{serving} (T-4).
\end{enumerate}

\section{Batasan Masalah}
\label{sec:batasan}

Batasan masalah pada Tugas Akhir ini menetapkan ruang lingkup yang sengaja tidak dicakup oleh kontribusi penelitian agar penilaian tetap tertuju pada arsitektur platform yang \textit{domain-agnostic}. Karena keluaran utama penelitian adalah arsitektur yang dapat dipakai lintas \textit{use case}, pilihan perwujudan seperti jumlah \textit{node}, patokan versi, atau sintaks komponen tertentu tidak dimasukkan sebagai batasan penelitian, melainkan dinyatakan sebagai batasan implementasi bersama lingkungan implementasi pada Bab~\ref{chap:implementasi}. Batasan-batasan di bawah ini membatasi klaim pada \textit{control plane} platform, cakupan keamanan antar layanan, dan lingkup evaluasi. Butir-butir tersebut selanjutnya dirujuk dengan singkatan BP-1 hingga BP-3, dan konsekuensinya pada evaluasi dibahas kembali pada Bab~\ref{chap:evaluasi}.

\vspace{0.5em}
\begin{enumerate}
    \item Fokus penelitian berada pada \textit{control plane} platform, bukan pengembangan model spesifik, sehingga \textit{use case} domain hanya dipakai sebagai kasus verifikasi rancangan dan pembahasan platform dijaga \textit{domain-agnostic}.

    \item Aspek keamanan berfokus pada akses antar layanan internal melalui \textit{mutual TLS} pada \textit{service mesh}, manajemen rahasia terpusat, dan kebijakan jaringan Kubernetes, sedangkan autentikasi pengguna akhir pada \textit{application layer} berada di luar cakupan penelitian.

    \item Evaluasi difokuskan pada penerimaan fungsional dan struktural tanpa mencakup \textit{deployment} produksi pada lingkungan multi-\textit{tenant} dengan trafik publik. Pengukuran latensi dan \textit{throughput} pada Bab~\ref{chap:evaluasi} bersifat indikatif pada lingkungan \textit{node} tunggal, sedangkan karakterisasi beban kuantitatif penuh pada skala produksi, seperti titik jenuh \textit{throughput} dan verifikasi skalabilitas horizontal lintas \textit{node}, ditempatkan sebagai arah pengembangan lanjutan.
\end{enumerate}

Konteks pembacaan seluruh laporan, terutama pada Bab~\ref{chap:analisis-masalah} dan Bab~\ref{chap:evaluasi}, berpijak pada ketiga batasan penelitian di atas. Batasan-batasan itu menjaga fokus kontribusi pada \textit{control plane} yang \textit{domain-agnostic} sehingga keberlakuan rancangan tidak terikat pada satu domain maupun satu lingkungan perwujudan. Pilihan perwujudan yang dapat diganti tanpa mengubah kontribusi, yaitu lingkungan \textit{k3s} satu \textit{node} dan patokan versi April 2026, dinyatakan sebagai batasan implementasi BI-1 dan BI-2 pada Bab~\ref{chap:implementasi}. Dengan pemisahan ini, keterbatasan lingkungan pengujian tidak membatasi keberlakuan arsitektur yang dirancang.

\section{Metodologi}
\label{sec:metodologi}

Penelitian ini mengikuti kerangka \textit{Design Science Research Methodology} (DSRM) yang dirumuskan oleh \textcite{peffers2007dsrm} untuk penelitian sistem informasi yang menghasilkan artefak. Keluaran utama Tugas Akhir ini berupa platform yang berjalan, yaitu artefak jenis instansiasi (\textit{instantiation}) dalam taksonomi \textit{design science}, sehingga kerangka penelitian berorientasi artefak menjadi pilihan yang tepat \autocite{hevner2004design}. DSRM dipilih dibandingkan metodologi rekayasa perangkat lunak murni karena DSRM menempatkan iterasi antara perancangan, demonstrasi, dan evaluasi sebagai bagian inti proses penelitian, bukan sekadar tahapan pengembangan. Kerangka proses \textit{data science} seperti CRISP-DM dan TDSP tidak menjadi pilihan karena keduanya menata pembangunan satu model pada satu domain data, sedangkan keluaran Tugas Akhir ini adalah platform lintas domain dan model hanya menjadi salah satu beban kerja yang dilayani. Tahapan yang dijalankan meliputi enam fase berikut dengan keluaran teknis yang dapat diukur.

\vspace{0.5em}
\begin{enumerate}
    \item \textbf{Identifikasi masalah dan motivasi}

    Studi literatur DataOps, MLOps, dan tata kelola data dilakukan untuk menyusun peta tantangan. Fase ini menghasilkan tiga pemicu dan empat rumusan masalah (RM-1 sampai RM-4).

    \vspace{0.5em}

    \item \textbf{Penetapan tujuan artefak}

    Tujuan dirumuskan sebagai daftar artefak teknis yang dapat dipasang, dioperasikan, dan diaudit, dipetakan satu lawan satu terhadap rumusan masalah. Kebutuhan fungsional KF-01 sampai KF-14 dan kebutuhan nonfungsional KNF-01 sampai KNF-12 disusun pada Bab~\ref{chap:analisis-masalah}.

    \vspace{0.5em}

    \item \textbf{Perancangan dan pengembangan artefak}

    Himpunan layanan representatif diklasifikasikan dari praktik industri menjadi sembilan \textit{layer} platform, komponen \textit{open source} dipilih untuk tiap \textit{layer}, lalu arsitektur kesembilan \textit{layer} tersebut dirancang sebagai \textit{control plane} deklaratif berbasis Kubernetes, bertolak dari alternatif arsitektur terpilih hasil matriks evaluasi multi-kriteria pada Bab~\ref{chap:analisis-masalah}. Pengembangan manifes dilakukan modular pada repositori Gitea dengan \textit{folder} per komponen dan \textit{overlay} Kustomize.

    \vspace{0.5em}

    \item \textbf{Demonstrasi}

    Platform dipasang pada lingkungan \textit{k3s} dan diberi beban kerja dari satu \textit{use case} domain nyata sebagai kasus verifikasi, yang dirinci pada Bab~\ref{chap:evaluasi}. Demonstrasi diatur sebagai empat tahap iteratif, dengan setiap tahap memvalidasi satu \textit{sub-sistem} sesuai T-1 hingga T-4.

    \vspace{0.5em}

    \item \textbf{Evaluasi}

    Evaluasi dilaksanakan pada akhir siklus implementasi dengan tiga dimensi, yaitu fungsional, nonfungsional (latensi indikatif, skalabilitas struktural, ketersediaan, observabilitas, keamanan, biaya), serta tata kelola (\textit{lineage} tertelusur, kontrak data ditegakkan, kebijakan akses dapat diaudit). Angka kuantitatif evaluasi dilaporkan pada Bab~\ref{chap:evaluasi}.

    \vspace{0.5em}

    \item \textbf{Komunikasi}

    Hasil rancangan dan implementasi didokumentasikan pada laporan ini, beserta repositori Gitea yang menjadi acuan reproduksi. Keluaran fase ini berupa laporan, repositori sumber, dan presentasi ujian akhir.
\end{enumerate}

Rangkaian enam fase di atas dijalankan secara iteratif, dengan kemungkinan kembali ke fase sebelumnya ketika temuan baru muncul pada fase berikutnya. Iterasi tersebut sejalan dengan karakter DSRM sebagai metodologi berorientasi artefak yang menempatkan umpan balik dari demonstrasi dan evaluasi sebagai pemicu perbaikan rancangan. Pola pengembangan serupa, yaitu membangun platform atau \textit{pipeline} MLOps berbasis Kubernetes lalu menilainya melalui demonstrasi dan eksperimen, dipakai oleh \textcite{hsu2024endtoend} dan \textcite{rodrigues2025architecture}. Pelaporan tiap fase disusun pada bab yang sesuai sehingga keterhubungan antara fase metodologi dan struktur laporan tetap eksplisit.

\section{Sistematika Penulisan}
\label{sec:sistematika}

Laporan Tugas Akhir disusun atas tujuh bab dengan urutan yang mengikuti DSRM dan praktik pelaporan rekayasa perangkat lunak. Pemetaan antara rumusan masalah, tujuan, \textit{sub-sistem}, dan butir kesimpulan dijaga konsisten lewat penomoran RM-N, T-N, dan butir pada Bab~\ref{chap:perancangan} hingga Bab~\ref{chap:penutup}. Bab IV dan Bab V membentuk pasangan perancangan dan implementasi yang dibaca berurutan, sedangkan Bab VI dan Bab VII menutup siklus dengan evaluasi serta kesimpulan dan saran. Urutan bab beserta isi utamanya dirinci pada poin-poin berikut.

\vspace{0.5em}
\begin{enumerate}
    \item \textbf{Bab I Pendahuluan}

    Memuat latar belakang, rumusan masalah, tujuan, batasan masalah, metodologi penelitian, serta sistematika penulisan.

    \vspace{0.5em}

    \item \textbf{Bab II Studi Literatur}

    Meninjau konsep DataOps dan MLOps, Kubernetes, manajemen data, layanan fitur, \textit{model lifecycle}, tata kelola data, deteksi \textit{drift}, observabilitas, praktik GitOps, dan keamanan platform sebagai landasan teoretis, lalu mengidentifikasi kesenjangan pada penelitian terkait.

    \vspace{0.5em}

    \item \textbf{Bab III Analisis Masalah}

    Memetakan kondisi sistem dan menyusun kebutuhan fungsional KF-01 hingga KF-14 serta kebutuhan nonfungsional KNF-01 hingga KNF-12 sebagai turunan empat tujuan. Matriks evaluasi multi-kriteria membandingkan empat alternatif arsitektur, lalu posisi platform terhadap penelitian terkait dipetakan berdasarkan kebutuhan yang sama.

    \vspace{0.5em}

    \item \textbf{Bab IV Perancangan Arsitektur Platform DataOps dan MLOps}

    Menetapkan sembilan \textit{layer} platform representatif dari praktik industri, memilih komponen \textit{open source} untuk tiap \textit{layer}, lalu menjabarkan perancangan arsitektur kesembilan \textit{layer} dan empat \textit{sub-sistem} yang menjawab T-1 hingga T-4. Diagram arsitektur disajikan sebagai satu arsitektur layanan terintegrasi beserta rincian beralat per \textit{layer}.

    \vspace{0.5em}

    \item \textbf{Bab V Implementasi Arsitektur Platform DataOps dan MLOps}

    Memaparkan implementasi tiap \textit{sub-sistem} yang sejalan dengan Bab~\ref{chap:perancangan}, dengan praktik GitOps melalui Argo CD pada repositori Gitea sebagai sumber kebenaran tunggal.

    \vspace{0.5em}

    \item \textbf{Bab VI Evaluasi}

    Berisi metode evaluasi, hasil pengukuran, dan pembahasan terhadap kebutuhan fungsional dan nonfungsional pada Bab~\ref{chap:analisis-masalah}, serta perbandingan biaya infrastruktur dan tenaga kerja terhadap skenario layanan terkelola.

    \vspace{0.5em}

    \item \textbf{Bab VII Penutup}

    Menyimpulkan jawaban atas T-1 hingga T-4 dan memberikan saran pengembangan lanjutan untuk lingkungan multi-\textit{node} dan multi-domain.
\end{enumerate}

Setelah Bab VII, dokumen ditutup oleh Daftar Pustaka yang memuat seluruh referensi terkutip dan Lampiran yang memuat berkas pendukung. Daftar Pustaka memakai gaya \textit{Chicago author-date} (\texttt{biblatex-chicago}) dengan label Indonesia agar konsisten dengan tubuh laporan. Lampiran memuat parameter penilaian pemilihan alat \textit{open source} dan katalog perintah verifikasi \textit{use case}. Penomoran lampiran mengikuti abjad (Lampiran A dan B) sebagaimana konvensi laporan Tugas Akhir di STI ITB.
